\begin{figure}[p]
\centering
\includegraphics[width=\textwidth]{supplementary/curated/measured_transfer_geometry.pdf}
\caption{\textbf{The measured-response learning comparison is limited by mapped coverage and transfer geometry.} \textbf{A}, Number of mapped functional presynaptic partners for each of seven targets in one mouse. \textbf{B}, Actual four-route support for the scan selected by median mapped-input count, with identifier tie breaking; rows are input coordinates and columns selected routes. The selected supports are mostly single distal sites, rather than illustrative nested bands. \textbf{C}, Response-prediction NMSE for exact learning, treeless ridge and restricted dictionaries, with random and shuffled surrogate rows; dots are seven target means, diamonds and bars their means and 95\% target-bootstrap intervals, and the dotted line marks ridge. \textbf{D}, Common-checkpoint update reconstruction by the unrestricted fixed transfer profile, the restricted ancestry profile and oracle trialwise ancestry amplitudes. The unrestricted fixed profile is nearly sufficient because somatic error mainly scales a stable transfer vector; restricted rules underperform ridge. \textbf{E}, Mean split-half Spearman reliability in 1,000 independent calibration datasets versus the observed value, for 125 partner records; the two negative estimates are assigned zero reliable variance and the identity line is not fitted. Scans and splits are averaged within targets before inference; $n$ is seven targets in \textbf{C,D} and 125 partner--scan records in \textbf{E}, with 95\% bootstrap intervals throughout. Target-level associations for all four topology measures are in main Fig.~9C, and coverage across all thirteen scans is in main Fig.~9F. This is an offline passive-tree geometry diagnostic, distinct from the empirical ancestry-response association.}
\label{fig:si_measured_transfer_geometry}
\label{fig:supp_measured_detail}
\label{fig:supp_measured_alignment_power}
\label{fig:si_measured_topology}
\end{figure}

\begin{figure}[p]
\centering
\includegraphics[width=\textwidth]{supplementary/curated/measured_predictor_controls.pdf}
\caption{\textbf{Linear baselines explain most measured-response prediction, and degrading observed inputs worsens performance.} Thirteen scans retain the complete-tree analysis's 130 outer stimulus-identity splits. All preprocessing, reliability masks and ridge tuning use training identities only. \textbf{A}, Training-mean, ordinary least-squares (OLS) and nested ridge-regression baselines against the matched nonlinear fits. Lower normalized MSE is better. Below, paired exact-nonlinear-minus-ridge differences show the small additional nonlinear benefit: mean $-0.0152$ (95\% interval $-0.0271$ to $-0.0042$). \textbf{B}, Ridge prediction when retaining nested random fractions of observed presynaptic inputs; five masks per split are averaged before inference. Manual-only partners are a separate provenance restriction, not a random retention dose. Empty masks use the training mean and remain included. \textbf{C}, Added Gaussian predictor noise in both training and test data, in units of training input standard deviation (SD) after the quantile transform. \textbf{D}, Reliability filtering based only on training repeats; parentheses give mean retained partner counts. Thin traces or dots show seven target means after averaging splits within scan and scans within target. Larger symbols and bars are means and 95\% target-bootstrap intervals (20,000 draws). All targets come from one mouse. Only ridge is refitted under the sensitivity conditions. Removing observed inputs does not recover unknown partners, establish population power or explain away the absence of preferential ancestry alignment.}
\label{fig:si_measured_predictor_controls}
\label{fig:supp_response_baselines}
\end{figure}

\begin{figure}[p]
\centering
\includegraphics[width=\textwidth]{supplementary/curated/animal_credit_reanalysis.pdf}
\caption{\textbf{Published animal data are consistent with signed per-neuron teaching coordinates.} This reanalysis contains no reconstructed dendritic morphology and measures no within-tree learning signal; it tests only the sign of two neuron-level coordinates. \textbf{A}, In the source brain--computer-interface (BCI) experiment, P+ and P- populations had positive and negative causal decoder weights, predicting opposite dendritic residual contrasts. \textbf{B}, Descriptive neuron-level residual distributions from published source sheet ``5E'' for error-increasing (inc.) and error-reducing (red.) conditions, in the source z-score normalization; dots show means and bars show neuron-level standard errors. These bars are descriptive and are not animal-bootstrap intervals. The paired animal contrasts and mode-energy decomposition use $n=6$ animals and are reported in Table~\ref{tab:animal_credit}; animals are the inferential units for those tests. This retrospective external comparison is consistent with signed per-neuron information but does not resolve within-tree transport or uniquely identify a conductance mechanism. This reanalysis is restored as a figure so that the six-animal signed-credit comparison described in Methods has drawn support; the animal-level contrasts remain tabulated.}
\label{fig:si_animal_credit_reanalysis}
\label{fig:supp_animal_credit}
\end{figure}
